% !TEX root = ../../main.tex
\section{Mô hình Thực thể -- Liên kết mở rộng (ERD / EERD)}
\label{sec:1_3_mo_hinh_erd_eerd}

% =====================================================
% GIỚI THIỆU
% =====================================================
\subsection{Giới thiệu}
% Mục này mô tả mục đích của tài liệu: trình bày mô hình thực thể (ERD)
% mở rộng (EERD) cho hệ thống FamilyConnect, bao gồm các thực thể chính,
% thuộc tính khóa, thuộc tính số (kiểu dữ liệu) và các mối quan hệ giữa chúng.
Tài liệu này trình bày mô hình thực thể -- mối quan hệ (ERD) mở rộng (EERD)
cho nền tảng \textbf{FamilyConnect}. Mô hình bao gồm 11 thực thể chính,
thể hiện đầy đủ thuộc tính khóa (PK), thuộc tính khóa ngoại (FK), thuộc tính
số liệu, và các mối quan hệ (kèm bậc số -- cardinality) giữa các thực thể.

% =====================================================
% DANH SÁCH THỰC THỂ
% =====================================================
\subsection{Danh sách thực thể và thuộc tính}

% ---- USER ----
\subsubsection{USER (Người dùng)}
% Thực thể đại diện cho tài khoản đăng nhập hệ thống, phục vụ nhóm
% chức năng "Người dùng & Bảo mật" (đăng ký, xác thực, RBAC).
\begin{longtable}{@{}p{3cm}p{2cm}p{9.5cm}@{}}
\toprule
\textbf{Thuộc tính} & \textbf{Khóa} & \textbf{Ghi chú} \\
\midrule
id & PK & Định danh duy nhất (uuid) \\
username & & Tên đăng nhập \\
email & & Địa chỉ email, dùng xác thực \\
password\_hash & & Mật khẩu đã mã hóa \\
role & & Vai trò (RBAC): quản trị viên / thành viên \\
\bottomrule
\end{longtable}

% ---- FAMILY ----
\subsubsection{FAMILY (Gia đình)}
% Thực thể gốc đại diện cho một gia đình / chi / nhánh, là điểm neo
% cho toàn bộ dữ liệu gia phả, cộng đồng, sự kiện và di sản.
\begin{longtable}{@{}p{3cm}p{2cm}p{9.5cm}@{}}
\toprule
\textbf{Thuộc tính} & \textbf{Khóa} & \textbf{Ghi chú} \\
\midrule
id & PK & Định danh duy nhất \\
name & & Tên gia đình / dòng họ \\
branch & & Chi / nhánh / phái \\
description & & Mô tả tổng quan \\
\bottomrule
\end{longtable}

% ---- PERSON ----
\subsubsection{PERSON (Thành viên gia phả)}
% Thực thể lõi của phân hệ gia phả. Thuộc tính parent_id tự tham chiếu
% (self-reference) đến chính thực thể PERSON để biểu diễn quan hệ
% cha/mẹ - con qua nhiều thế hệ mà không cần bảng riêng cho từng cấp.
\begin{longtable}{@{}p{3cm}p{2cm}p{9.5cm}@{}}
\toprule
\textbf{Thuộc tính} & \textbf{Khóa} & \textbf{Ghi chú} \\
\midrule
id & PK & Định danh duy nhất \\
family\_id & FK & Tham chiếu đến FAMILY \\
user\_id & FK & Tham chiếu đến USER (nếu thành viên có tài khoản) \\
parent\_id & FK & Tự tham chiếu đến PERSON (quan hệ cha/mẹ) \\
full\_name & & Họ và tên đầy đủ \\
dob & & Ngày sinh \\
dod & & Ngày mất (nếu có) \\
gender & & Giới tính \\
occupation & & Nghề nghiệp / học vấn \\
\bottomrule
\end{longtable}

% ---- MARRIAGE ----
\subsubsection{MARRIAGE (Hôn nhân)}
% Tách thành bảng riêng thay vì gắn trực tiếp vào PERSON, vì một người
% có thể có nhiều cuộc hôn nhân theo thời gian (đảm bảo chuẩn hóa dữ liệu).
\begin{longtable}{@{}p{3cm}p{2cm}p{9.5cm}@{}}
\toprule
\textbf{Thuộc tính} & \textbf{Khóa} & \textbf{Ghi chú} \\
\midrule
id & PK & Định danh duy nhất \\
person1\_id & FK & Tham chiếu đến PERSON (vợ/chồng 1) \\
person2\_id & FK & Tham chiếu đến PERSON (vợ/chồng 2) \\
marriage\_date & & Ngày kết hôn \\
\bottomrule
\end{longtable}

% ---- POST / COMMENT ----
\subsubsection{POST (Bài đăng) \& COMMENT (Bình luận)}
% Phục vụ nhóm chức năng "Cộng đồng & Tương tác": chia sẻ tin tức,
% hình ảnh gia đình và tương tác bình luận.
\begin{longtable}{@{}p{3cm}p{2cm}p{9.5cm}@{}}
\toprule
\textbf{Thuộc tính} & \textbf{Khóa} & \textbf{Ghi chú} \\
\midrule
\multicolumn{3}{l}{\textit{POST}} \\
id & PK & Định danh duy nhất \\
user\_id & FK & Người đăng bài \\
family\_id & FK & Bài đăng thuộc gia đình nào \\
content & & Nội dung bài đăng \\
created\_at & & Thời điểm đăng \\
\addlinespace
\multicolumn{3}{l}{\textit{COMMENT}} \\
id & PK & Định danh duy nhất \\
post\_id & FK & Tham chiếu đến POST \\
user\_id & FK & Người bình luận \\
content & & Nội dung bình luận \\
created\_at & & Thời điểm bình luận \\
\bottomrule
\end{longtable}

% ---- EVENT / RSVP ----
\subsubsection{EVENT (Sự kiện) \& RSVP}
% Phục vụ nhóm "Quản lý Sự kiện": tạo sự kiện, nhắc nhở, quản lý người
% tham gia. RSVP dùng khóa chính kép (event_id, user_id) -- không cần id riêng.
\begin{longtable}{@{}p{3cm}p{2cm}p{9.5cm}@{}}
\toprule
\textbf{Thuộc tính} & \textbf{Khóa} & \textbf{Ghi chú} \\
\midrule
\multicolumn{3}{l}{\textit{EVENT}} \\
id & PK & Định danh duy nhất \\
family\_id & FK & Sự kiện thuộc gia đình nào \\
name & & Tên sự kiện \\
event\_date & & Ngày diễn ra \\
location & & Địa điểm \\
\addlinespace
\multicolumn{3}{l}{\textit{RSVP}} \\
event\_id & PK, FK & Khóa chính kép -- Tham chiếu EVENT \\
user\_id & PK, FK & Khóa chính kép -- Tham chiếu USER \\
status & & Trạng thái tham gia (đồng ý / từ chối / chưa rõ) \\
\bottomrule
\end{longtable}

% ---- DOCUMENT ----
\subsubsection{DOCUMENT (Tài liệu / Di sản)}
% Phục vụ nhóm "Danh bạ & Di sản": lưu trữ tài liệu, câu chuyện lịch sử,
% lưu trữ số của gia đình.
\begin{longtable}{@{}p{3cm}p{2cm}p{9.5cm}@{}}
\toprule
\textbf{Thuộc tính} & \textbf{Khóa} & \textbf{Ghi chú} \\
\midrule
id & PK & Định danh duy nhất \\
family\_id & FK & Tài liệu thuộc gia đình nào \\
title & & Tiêu đề tài liệu \\
file\_url & & Đường dẫn lưu trữ tệp \\
story\_text & & Nội dung câu chuyện lịch sử \\
\bottomrule
\end{longtable}

% ---- NOTIFICATION ----
\subsubsection{NOTIFICATION (Thông báo)}
% Phục vụ hệ thống thông báo trong nhóm "Cộng đồng & Tương tác".
\begin{longtable}{@{}p{3cm}p{2cm}p{9.5cm}@{}}
\toprule
\textbf{Thuộc tính} & \textbf{Khóa} & \textbf{Ghi chú} \\
\midrule
id & PK & Định danh duy nhất \\
user\_id & FK & Người nhận thông báo \\
content & & Nội dung thông báo \\
is\_read & & Trạng thái đã đọc / chưa đọc \\
\bottomrule
\end{longtable}

% =====================================================
% MỐI QUAN HỆ
% =====================================================
\subsection{Danh sách mối quan hệ (Relationships)}
% Bảng tổng hợp các mối quan hệ giữa các thực thể, kèm bậc số
% (cardinality) theo ký hiệu chuẩn ERD: 1 - n (một - nhiều).
\begin{longtable}{@{}p{3.2cm}p{1.2cm}p{3.2cm}p{6.5cm}@{}}
\toprule
\textbf{Thực thể A} & \textbf{Bậc số} & \textbf{Thực thể B} & \textbf{Diễn giải} \\
\midrule
USER & 1--n & PERSON & Một tài khoản có thể liên kết một hồ sơ gia phả \\
FAMILY & 1--n & PERSON & Một gia đình gồm nhiều thành viên \\
PERSON & 1--n & PERSON & Quan hệ cha/mẹ -- con (tự tham chiếu) \\
PERSON & 1--n & MARRIAGE & Một người có thể có nhiều cuộc hôn nhân \\
USER & 1--n & POST & Một người dùng đăng nhiều bài viết \\
FAMILY & 1--n & POST & Bài viết thuộc về một gia đình \\
POST & 1--n & COMMENT & Một bài viết nhận nhiều bình luận \\
USER & 1--n & COMMENT & Một người dùng viết nhiều bình luận \\
FAMILY & 1--n & EVENT & Một gia đình tổ chức nhiều sự kiện \\
EVENT & 1--n & RSVP & Một sự kiện nhận nhiều phản hồi tham gia \\
USER & 1--n & RSVP & Một người dùng tham gia nhiều sự kiện \\
FAMILY & 1--n & DOCUMENT & Một gia đình lưu trữ nhiều tài liệu \\
USER & 1--n & NOTIFICATION & Một người dùng nhận nhiều thông báo \\
\bottomrule
\end{longtable}
